% problem_id: S001-theorem-normalized-base-change-with-reduced-fibre
% source_id: S001 (Stacks Project)
% source_file: more-morphisms.tex
% statement_locator: more-morphisms.tex lines 19281--19289
% proof_locator: more-morphisms.tex lines 19291--19355
% env: theorem
% id_note: label 在本来源内唯一
% status: complete (statement + author proof verbatim + reason + locators)
%
\begin{theorem}
\label{theorem-normalized-base-change-with-reduced-fibre}
Let $A$ be a Dedekind ring with fraction field $K$.
Let $X$ be a scheme flat and of finite type over $A$.
Assume $A$ is a Nagata ring.
There exists a finite extension $L/K$ such that
the normalized base change $Y$ is smooth over $\Spec(B)$
at all generic points of all fibres.
\end{theorem}

\begin{proof}
During the proof we will repeatedly use that formation of the set of points
where a (flat, finitely presented) morphism like $X \to \Spec(A)$ is
smooth commutes with base change, see
Morphisms, Lemma \ref{morphisms-lemma-set-points-where-fibres-smooth}.

\medskip\noindent
We first choose a finite extension $L/K$ such that
$(X_L)_{red}$ is geometrically reduced over $L$, see
Varieties, Lemma \ref{varieties-lemma-finite-extension-geometrically-reduced}.
Since $Y \to (X_B)_{red}$ is birational we see applying
Varieties, Lemma \ref{varieties-lemma-generic-points-geometrically-reduced}
that $Y_L$ is geometrically reduced over $L$ as well.
Hence $Y_L \to \Spec(L)$ is smooth on a dense open $V \subset Y_L$ by
Varieties, Lemma \ref{varieties-lemma-geometrically-reduced-dense-smooth-open}.
Thus the smooth locus $U \subset Y$ of the morphism $Y \to \Spec(B)$
is open (by Morphisms, Definition \ref{morphisms-definition-smooth})
and is dense in the generic fibre. Replacing $A$ by $B$ and $X$ by $Y$
we reduce to the case treated in the next paragraph.

\medskip\noindent
Assume $X$ is normal and the smooth locus $U \subset X$ of $X \to \Spec(A)$
is dense in the generic fibre. This implies that $U$ is dense in all but
finitely many fibres, see Lemma \ref{lemma-nowhere-dense-generic-fibre}.
Let $x_1, \ldots, x_r \in X \setminus U$ be the finitely many generic
points of irreducible components of $X \setminus U$ which are moreover
generic points of irreducible components of fibres of $X \to \Spec(A)$.
Set $\mathcal{O}_i = \mathcal{O}_{X, x_i}$. Let $A_i$ be the localization of
$A$ at the maximal ideal corresponding to the image of $x_i$ in $\Spec(A)$.
By
More on Algebra, Proposition
\ref{more-algebra-proposition-epp-essentially-finite-type}
there exist finite extensions
$K_i/K$ which are solutions for the extension of discrete valuation
rings $A_i \to \mathcal{O}_i$. Let $L/K$ be a finite extension
dominating all of the extensions $K_i/K$. Then $L/K$
is still a solution for $A_i \to \mathcal{O}_i$ by
More on Algebra, Lemma \ref{more-algebra-lemma-solution-goes-up}.

\medskip\noindent
Consider the diagram (\ref{equation-normalized-base-change})
with the extension $L/K$ we just produced. Note that $U_B \subset X_B$
is smooth over $B$, hence normal (for example use
Algebra, Lemma \ref{algebra-lemma-normal-goes-up}).
Thus $Y \to X_B$ is an isomorphism over $U_B$.
Let $y \in Y$ be a generic point of an irreducible
component of a fibre of $Y \to \Spec(B)$ lying over the maximal
ideal $\mathfrak m \subset B$. Assume that $y \not \in U_B$.
Then $y$ maps to one of the points $x_i$. It follows that
$\mathcal{O}_{Y, y}$ is a local ring of the integral closure
of $\mathcal{O}_i$ in $R(X) \otimes_K L$ (details omitted).
Hence because $L/K$ is a solution for
$A_i \to \mathcal{O}_i$ we see that
$B_\mathfrak m \to \mathcal{O}_{Y, y}$ is formally smooth
in the $\mathfrak m_y$-adic topology
(this is the definition of being a "solution").
In other words, $\mathfrak m\mathcal{O}_{Y, y} = \mathfrak m_y$
and the residue field extension is separable, see
More on Algebra, Lemma \ref{more-algebra-lemma-extension-dvrs-formally-smooth}.
Hence the local ring
of the fibre at $y$ is $\kappa(y)$.
This implies the fibre is smooth over $\kappa(\mathfrak m)$
at $y$ for example by Algebra, Lemma \ref{algebra-lemma-separable-smooth}.
This finishes the proof.
\end{proof}
